\section{Differential Geometry Preview: Calculus on Curved Worlds}
\chapterline

\subsection{The geodetic surveyor: how mountain peaks forced a new geometry}

In Chapter~22, we set aside the ruler. By discarding distances and angles, topology revealed the qualitative structure of space---open sets, connected components, and boundary-free loops that survive continuous stretching.

Yet physics and engineering cannot live on rubber sheets alone. To calculate planetary orbits, predict fluid flow across a curved wing, or describe gravitational fields, we need quantitative calculus. We must measure velocities, compute accelerations, evaluate line integrals, and solve differential equations.

The dilemma is immediate. If you stand in the middle of an open field and look around, the ground appears flat. You can lay down a Cartesian coordinate grid, set up perpendicular $x$ and $y$ axes, and use standard multivariable calculus to compute gradients, tangent lines, and areas.

Yet if you zoom out into space, you see that you are standing on the surface of a sphere. The Earth has no global $xy$-plane. Any attempt to flatten the entire globe onto a single sheet of paper tears continents or distorts distances beyond recognition.

Between 1821 and 1825, King George IV commissioned Carl Friedrich Gauss to lead the \emph{geodetic survey} of the Kingdom of Hanover, measuring the terrain and computing accurate distances across the region. During the survey, Gauss used the \emph{heliotrope} (a precision mirror device to reflect flashes of sunlight across dozens of kilometers) to measure large triangles between the peaks of the Brocken, the Hoher Hagen, and the Inselsberg.

Working with these large-scale measurements led Gauss to consider whether an observer living \emph{entirely within} a two-dimensional surface could determine whether the surface was flat or curved, without referencing an ambient third dimension.

In his 1827 treatise, \emph{Disquisitiones generales circa superficies curvas} (``General Investigations of Curved Surfaces''), Gauss proved his \textbf{Theorema Egregium} (the ``Remarkable Theorem''): \emph{curvature is intrinsic}. An observer living on a surface can measure its curvature purely by taking distance and angle measurements along paths within the surface itself, without any reference to an ambient higher-dimensional space. By contrast, \textbf{extrinsic curvature} is bending that exists only relative to an external surrounding space. Rolling a flat sheet of paper into a cylinder curves it in three dimensions, but an ant living on the paper still measures standard flat Euclidean distances: its intrinsic curvature is zero. On a sphere, however, the curvature is genuinely intrinsic: no patch of a sphere can be flattened without stretching or tearing.

\subsection{The challenge of curved space}

This is the central idea of \textbf{differential geometry}: doing calculus on curved spaces by approximating them locally with flat linear spaces. Because curved spaces rarely admit a single global Cartesian grid, derivatives, velocity vectors, and integrals have to be built out of local flat patches.

The strategy is built on three interconnected layers, illustrated in Figure~\ref{23:fig-chart-tangent}:
\begin{enumerate}
\item \textbf{The global curved space} (the manifold $M$).
\item \textbf{A local flat translation} (the coordinate chart $\phi$), which maps an open patch $U \subset M$ into an ordinary open subset of $\R^n$.
\item \textbf{A local linear space} (the tangent plane $T_p M$), where instantaneous velocities and directional derivatives live as true vectors.
\end{enumerate}

\begin{figure}[htbp]
\centering
\begin{tikzpicture}[scale=1.0, >=Latex]
  % 3D Curved Surface M
  \draw[thick, sourceblue, fill=sourceblue!10] 
    (-4.5, -0.5) to[out=30, in=150] (-0.5, 0.5) 
    to[out=-30, in=120] (0.5, -2.5) 
    to[out=-160, in=-20] (-3.5, -2.2) 
    to[out=160, in=-60] (-4.5, -0.5);
  \node[font=\bfseries, text=sourceblue, anchor=south east] at (-2.7, 0.38) {Manifold $M$};

  % Open patch U on surface
  \draw[thick, dashed, targetgreen, fill=targetgreen!15]
    (-2.8, -0.7) to[out=20, in=160] (-1.2, -0.5)
    to[out=-30, in=100] (-1.0, -1.8)
    to[out=-160, in=-20] (-2.6, -1.7)
    to[out=150, in=-70] (-2.8, -0.7);
  \node[font=\footnotesize\bfseries, text=targetgreen, anchor=south] at (-2.0, -0.42) {Patch $U$};

  % Tangent plane at p (parallelogram - draw behind point and curve)
  \draw[thick, quotientorange, fill=quotientorange!20, fill opacity=0.7]
    (-3.0, -1.8) -- (-0.8, -1.5) -- (-0.4, -0.3) -- (-2.6, -0.6) -- cycle;
  \node[font=\footnotesize\bfseries, text=quotientorange, right=4pt] at (-0.5, -1.1) {Tangent Plane $T_p M$};

  % Point p coordinates
  \coordinate (P) at (-1.8, -1.1);

  % Curve gamma(t) on surface
  \draw[thick, structurepurple] (-2.5, -1.5) to[out=20, in=200] (P) to[out=20, in=180] (-1.1, -0.8);
  \node[font=\scriptsize, text=structurepurple, anchor=north] at (-2.65, -1.85) {Curve $\gamma(t)$};

  % Point p (drawn on top of plane and curve)
  \fill[structurepurple] (P) circle (2.2pt);
  \node[below left=2pt, font=\small\bfseries, text=structurepurple] at (P) {$p$};

  % Tangent vector v
  \draw[->, very thick, trapred] (P) -- ++(35:1.2);
  \node[trapred, font=\footnotesize\bfseries, fill=white, inner sep=1pt, anchor=south west] at (-0.78, -0.22) {$v = \gamma'(0)$};

  % Downward chart map arrow phi
  \draw[->, very thick, targetgreen] (-1.9, -2.6) -- (-1.9, -3.8) 
    node[midway, right=4pt, font=\small\bfseries] {Chart map $\phi$};

  % Flat Euclidean space R^2 below
  \draw[thick, gray!60, fill=gray!8, rounded corners=4pt] (-4.0, -6.0) rectangle (0.2, -4.0);
  \draw[->, gray!70] (-3.8, -5.8) -- (-3.8, -4.2) node[left=2pt, font=\scriptsize] {$u_2$};
  \draw[->, gray!70] (-3.8, -5.8) -- (-0.6, -5.8) node[below=2pt, font=\scriptsize] {$u_1$};

  % Flat image phi(U)
  \draw[thick, targetgreen, fill=targetgreen!25] 
    (-2.8, -5.5) to[out=30, in=170] (-1.0, -5.4)
    to[out=-20, in=100] (-0.8, -4.6)
    to[out=-160, in=20] (-2.4, -4.7)
    to[out=160, in=-80] (-2.8, -5.5);
  \node[font=\small, text=targetgreen] at (-1.8, -4.3) {Flat domain $\phi(U) \subset \R^2$};

  \coordinate (PhiP) at (-1.8, -5.1);
  \fill[structurepurple] (PhiP) circle (2pt);
  \node[below=3pt, font=\small\bfseries, text=structurepurple] at (PhiP) {$\phi(p)$};
\end{tikzpicture}
\caption{The local architecture of a smooth manifold. A curved manifold $M$ has no single global Cartesian coordinate grid. Instead, a coordinate chart $\phi$ maps a local open patch $U \subset M$ homeomorphically to a flat domain $\phi(U) \subset \R^n$. Directional derivatives along curves $\gamma(t)$ through $p$ form the tangent space $T_p M$, a flat vector space touching $M$ at $p$.}\label{23:fig-chart-tangent}
\end{figure}

\subsection{Manifolds: locally Euclidean worlds}

The central concept that reconciles local flatness with global curvature is the \textbf{manifold}: a space that may bend and twist globally, but looks indistinguishable from flat Euclidean space to anyone examining a tiny patch around themselves.

\begin{definition}
An \textbf{$n$-dimensional topological manifold} $M$ is a topological space where every point $p \in M$ possesses an open neighborhood $U \subseteq M$ homeomorphic to an open subset of $\R^n$.
\end{definition}

Recall from Chapter~22 that a \textbf{homeomorphism} is a continuous bijection whose inverse is also continuous: a rubber-sheet equivalence that stretches and bends space freely without ever tearing or gluing. This means every point on an $n$-dimensional manifold sits inside a small neighborhood that is topologically identical to a flat patch of $\R^n$.
\begin{itemize}
\item A 1-dimensional manifold locally resembles a line (such as a circle $S^1$ or a smooth wire).
\item A 2-dimensional manifold locally resembles the plane (such as a sphere $S^2$, a torus, or a saddle).
\item A 4-dimensional manifold locally resembles $\R^4$ (such as spacetime in general relativity).
\end{itemize}

\subsection{Atlases and smooth coordinate charts}

Topological equivalence alone is not enough to do calculus. A rubber-sheet homeomorphism allows sharp corners and crumpled bends where derivatives do not exist. To compute velocities and gradients, we need to differentiate.

A single flat map cannot cover a sphere without pinching, tearing, or creating singular points. Yet a collection of overlapping flat maps can cover the whole globe cleanly. In cartography, a bound book of regional flat maps that together describe the entire Earth is called an \textbf{atlas}. Differential geometry adopts this exact terminology:
\begin{itemize}
\item A \textbf{coordinate chart} is a pair $(U, \phi)$, where $U \subseteq M$ is an open patch and $\phi: U \to \R^n$ is a homeomorphism assigning concrete numerical coordinates $(x_1, \dots, x_n)$ to points in $U$.
\item An \textbf{atlas} is a collection of charts $\{(U_\alpha, \phi_\alpha)\}$ whose domains together cover all of $M$.
\end{itemize}

Whenever two charts $(U_\alpha, \phi_\alpha)$ and $(U_\beta, \phi_\beta)$ overlap, any point in the intersection $U_\alpha \cap U_\beta$ receives two competing coordinate descriptions. The dictionary converting one coordinate system into the other is the \textbf{transition map}:
\[
\tau_{\alpha \beta} = \phi_\beta \circ \phi_\alpha^{-1} : \phi_\alpha(U_\alpha \cap U_\beta) \to \phi_\beta(U_\alpha \cap U_\beta).
\]
In this construction (illustrated in Figure~\ref{23:fig-transition-map}): $\phi_\alpha^{-1}$ lifts flat coordinates from $\R^n$ back onto the curved manifold patch, and $\phi_\beta$ maps that curved patch down into a different region of flat $\R^n$. Therefore, the transition map $\tau_{\alpha \beta}$ is an ordinary function from an open subset of $\R^n$ to $\R^n$!

\begin{figure}[htbp]
\centering
\begin{tikzpicture}[scale=0.92, >=Latex]
  % Top: Manifold M with overlapping patches (fills first)
  \draw[thick, gray!40, fill=gray!6] 
    (-4.8, 1.4) to[out=25, in=155] (0, 2.3) 
    to[out=-25, in=155] (4.8, 1.6) 
    to[out=-60, in=30] (4.0, 0.1) 
    to[out=-150, in=-30] (-4.0, 0.0) 
    to[out=150, in=-70] cycle;

  % Open patch U_alpha fill & border
  \draw[thick, sourceblue, dashed, fill=sourceblue!15] (-1.1, 1.1) ellipse (1.6 and 0.85);

  % Open patch U_beta fill & border
  \draw[thick, targetgreen, dashed, fill=targetgreen!15] (1.1, 1.1) ellipse (1.6 and 0.85);

  % Overlap intersection fill
  \begin{scope}
    \clip (-1.1, 1.1) ellipse (1.6 and 0.85);
    \fill[quotientorange!35] (1.1, 1.1) ellipse (1.6 and 0.85);
  \end{scope}

  % Manifold labels and points on top of all fills
  \node[font=\bfseries, text=gray!70!black] at (-3.8, 2.1) {Manifold $M$};
  \node[font=\footnotesize\bfseries, text=sourceblue] at (-2.3, 1.8) {Patch $U_\alpha$};
  \node[font=\footnotesize\bfseries, text=targetgreen] at (2.3, 1.8) {Patch $U_\beta$};
  \node[font=\scriptsize\bfseries, text=quotientorange!90!black] at (0, 1.8) {$U_\alpha \cap U_\beta$};

  % Point p in the overlap
  \fill[structurepurple] (0, 1.05) circle (2.2pt);
  \node[right=2pt, font=\small\bfseries, text=structurepurple] at (0, 1.05) {$p$};

  % Bottom Left: Flat chart space R^n (Chart alpha)
  \draw[thick, gray!50, fill=gray!5, rounded corners=4pt] (-5.8, -3.6) rectangle (-1.0, -1.0);
  \node[anchor=north west, font=\scriptsize, text=gray!70!black] at (-5.6, -1.1) {Flat domain $\R^n$};

  % Image phi_alpha(U_alpha)
  \draw[thick, sourceblue, fill=sourceblue!15] (-3.4, -2.3) ellipse (1.5 and 0.75);
  \node[font=\scriptsize\bfseries, text=sourceblue] at (-4.2, -2.3) {$\phi_\alpha(U_\alpha)$};

  % Image of overlap in chart alpha
  \draw[thick, quotientorange, fill=quotientorange!35] (-2.5, -2.3) ellipse (0.45 and 0.55);
  \fill[structurepurple] (-2.5, -2.3) circle (2pt);
  \node[above=2pt, font=\scriptsize\bfseries, text=structurepurple] at (-2.5, -2.3) {$\phi_\alpha(p)$};

  % Bottom Right: Flat chart space R^n (Chart beta)
  \draw[thick, gray!50, fill=gray!5, rounded corners=4pt] (1.0, -3.6) rectangle (5.8, -1.0);
  \node[anchor=north east, font=\scriptsize, text=gray!70!black] at (5.6, -1.1) {Flat domain $\R^n$};

  % Image phi_beta(U_beta)
  \draw[thick, targetgreen, fill=targetgreen!15] (3.4, -2.3) ellipse (1.5 and 0.75);
  \node[font=\scriptsize\bfseries, text=targetgreen] at (4.2, -2.3) {$\phi_\beta(U_\beta)$};

  % Image of overlap in chart beta
  \draw[thick, quotientorange, fill=quotientorange!35] (2.5, -2.3) ellipse (0.45 and 0.55);
  \fill[structurepurple] (2.5, -2.3) circle (2pt);
  \node[above=2pt, font=\scriptsize\bfseries, text=structurepurple] at (2.5, -2.3) {$\phi_\beta(p)$};

  % Downward chart mapping arrows
  \draw[->, very thick, sourceblue] (-2.2, 0.2) to[out=-130, in=90] 
    node[midway, left=7pt, font=\small\bfseries] {$\phi_\alpha$} (-3.4, -0.9);
  \draw[->, very thick, targetgreen] (2.2, 0.2) to[out=-50, in=90] 
    node[midway, right=7pt, font=\small\bfseries] {$\phi_\beta$} (3.4, -0.9);

  % Transition map arrow between the two flat overlap images
  \draw[->, very thick, quotientorange] (-1.9, -2.3) -- (1.9, -2.3)
    node[midway, above=6pt, font=\small\bfseries, text=quotientorange] {$\tau_{\alpha\beta} = \phi_\beta \circ \phi_\alpha^{-1}$};
  \node[font=\scriptsize, text=quotientorange!90!black, align=center] at (0, -3.0) {
    Smooth transition map between flat subsets of $\R^n$\\[1pt]
    \textcolor{gray!80!black}{(infinitely differentiable $C^\infty$)}
  };
\end{tikzpicture}
\caption{The transition map between overlapping charts on a smooth manifold. When open patches $U_\alpha$ and $U_\beta$ overlap on $M$, any point $p \in U_\alpha \cap U_\beta$ maps to both $\phi_\alpha(p)$ and $\phi_\beta(p)$ in flat Euclidean space $\R^n$. The transition map $\tau_{\alpha\beta} = \phi_\beta \circ \phi_\alpha^{-1}$ translates directly between these two Euclidean representations. A manifold is smooth if all its transition maps are smooth ($C^\infty$) in the standard multivariable calculus sense.}\label{23:fig-transition-map}
\end{figure}

If all transition maps in an atlas are infinitely differentiable ($C^\infty$), the atlas is called smooth, and $M$ is a \textbf{smooth manifold}. This immediately solves our foundational problem: we can define what it means for a function $f: M \to \R$ to be smooth by pulling it back to local coordinate charts ($f \circ \phi^{-1}: \R^n \to \R$) and using standard multivariable calculus.

\begin{intuitionbox}
\textbf{Pause and think.} Can a single coordinate chart cover the entire circle $S^1$ without missing a point or tearing?

No. If a single chart covered all of $S^1$, there would exist a continuous bijection $\phi: S^1 \to V$ where $V \subset \R$ is an open set. But the circle $S^1$ is compact (closed and bounded in $\R^2$), whereas no non-empty open subset of $\R$ is compact. In addition, removing any single point from $S^1$ leaves it connected, whereas removing any point from an open interval splits it into two disconnected pieces. A closed loop cannot be unfolded onto a single open line without a seam. To map the circle smoothly, we must use an atlas of at least two overlapping charts.
\end{intuitionbox}

\subsection{Tangent spaces: local linear approximations}

In single-variable calculus, the derivative of a function at a point gives the slope of the tangent line. On a curved surface $M$, what plays the role of the tangent line?

Imagine a particle traveling along a smooth curve $\gamma: (-\eps, \eps) \to M$ lying on the surface, with $\gamma(0) = p$. At the instant $t = 0$, the particle has an instantaneous velocity vector $\gamma'(0)$.

If you gather the velocity vectors of \emph{all} possible smooth curves passing through $p$, they fill out a flat vector space: the \textbf{tangent space} $T_p M$.
\begin{itemize}
\item If $M$ is a 2D surface, $T_p M$ is a 2D plane touching the surface at $p$.
\item The tangent space is a true vector space: you can add tangent vectors and scale them by real numbers, even though the curved surface itself has no vector addition!
\end{itemize}

Derivatives of maps between manifolds then become linear transformations between tangent spaces:
\[
df_p : T_p M \to T_{f(p)} N.
\]
Linear algebra provides the local description of non-linear geometry.

\subsection{The metric tensor: measuring lengths and angles}

A smooth manifold gives us coordinate charts and tangent spaces, but it does not yet give us metric geometry. A vector in $T_p M$ specifies a direction, but we cannot yet measure its length, nor can we calculate the angle between two intersecting curves or the area of a surface region.

To do metric geometry, we equip the manifold with a \textbf{Riemannian metric}: a smoothly varying inner product assigned to each tangent space $T_p M$.
\begin{definition}
A \textbf{Riemannian metric} $g$ on a smooth manifold $M$ assigns to each point $p \in M$ an inner product
\[
g_p : T_p M \times T_p M \to \R
\]
that is symmetric ($g_p(u, v) = g_p(v, u)$), positive-definite ($g_p(v, v) > 0$ for all $v \neq 0$), and depends smoothly on $p$.
\end{definition}

In a local coordinate chart $(x^1, \dots, x^n)$, the metric is represented by an $n \times n$ symmetric matrix of functions $g_{ij}(x) = g\left(\frac{\partial}{\partial x^i}, \frac{\partial}{\partial x^j}\right)$. The squared distance element is given by the \textbf{line element}:
\[
ds^2 = \sum_{i=1}^n \sum_{j=1}^n g_{ij}(x) \, dx^i \, dx^j.
\]
Using the metric tensor, we can compute the length of any smooth curve $\gamma: [a, b] \to M$ by integrating its instantaneous speed:
\[
L(\gamma) = \int_a^b \sqrt{g_{\gamma(t)}(\gamma'(t), \gamma'(t))} \, dt.
\]

\begin{example}[The metric of a sphere]
Consider a sphere of radius $R$. In spherical coordinates $(\theta, \phi)$, where $\theta \in (0, \pi)$ is colatitude (the angle measured down from the North Pole) and $\phi \in [0, 2\pi)$ is longitude, the line element is:
\[
ds^2 = R^2 \, d\theta^2 + R^2 \sin^2\theta \, d\phi^2.
\]
Here the metric matrix is diagonal: $g_{\theta\theta} = R^2$, $g_{\phi\phi} = R^2 \sin^2\theta$, and $g_{\theta\phi} = 0$.

Notice what this formula reveals about the geometry of the sphere without referencing an ambient three-dimensional space:
\begin{itemize}
\item Walking along a meridian (constant longitude, $d\phi = 0$), the distance element is $ds = R \, d\theta$. Every meridian behaves like a circle of radius $R$.
\item Walking along a parallel of latitude (constant $\theta$, $d\theta = 0$), the distance element is $ds = R \sin\theta \, d\phi$. Near the equator ($\theta = \pi/2$), $\sin\theta = 1$, giving a circumference of $2\pi R$. Near the North Pole ($\theta \to 0$), the factor $\sin^2\theta \to 0$, capturing how lines of longitude converge to a single point.
\end{itemize}
The functions $g_{ij}$ completely dictate the intrinsic geometry of the sphere.
\end{example}

\subsection{Geodesics, parallel transport, and curvature}

Once we can measure lengths, two natural questions arise: what is the straightest possible path between two points on a curved surface, and how does a directional vector change as you carry it along that path?

In flat Euclidean space, the shortest path between two points is a straight line, characterized by zero acceleration: $\gamma''(t) = 0$. On a curved manifold, the counterpart is a \textbf{geodesic}: a curve whose acceleration vector has zero tangential component along the surface. Along a geodesic, the curve bends only as much as the surface itself requires. On a sphere, the geodesics are great circles, such as the equator and the meridians.

Now imagine carrying an arrow (a tangent vector) while walking along a path. On flat paper, you keep the arrow pointing in the same direction simply by keeping its Cartesian components constant. But on a curved surface, the tangent plane tilts at every step. \textbf{Parallel transport} is the rule for sliding a vector along a curve while keeping its derivative along the surface equal to zero.

What happens when you parallel-transport a vector around a closed loop?

\begin{figure}[htbp]
\centering
\begin{tikzpicture}[scale=1.15, >=Latex]
  % Shaded sphere background
  \fill[sourceblue!6] (0, 0) circle (2.2);

  % Geodesic triangle NAB shaded
  \fill[quotientorange!20] (0, 2.2) -- (0, -0.65) 
    arc[start angle=270, end angle=360, x radius=2.2, y radius=0.65] 
    arc[start angle=0, end angle=90, radius=2.2] -- cycle;

  % Outer circle of sphere
  \draw[thick, sourceblue] (0, 0) circle (2.2);

  % Equator: front solid, back dashed
  \draw[thick, sourceblue] (-2.2, 0) arc[start angle=180, end angle=360, x radius=2.2, y radius=0.65];
  \draw[dashed, sourceblue!45] (2.2, 0) arc[start angle=0, end angle=180, x radius=2.2, y radius=0.65];

  % Prime Meridian NA (vertical in this projection)
  \draw[very thick, quotientorange!80!black] (0, 2.2) -- (0, -0.65);

  % Equator arc AB
  \draw[very thick, quotientorange!80!black] (0, -0.65) arc[start angle=270, end angle=360, x radius=2.2, y radius=0.65];

  % Meridian BN (circle arc)
  \draw[very thick, quotientorange!80!black] (2.2, 0) arc[start angle=0, end angle=90, radius=2.2];

  % Vertices N, A, B
  \coordinate (N) at (0, 2.2);
  \coordinate (A) at (0, -0.65);
  \coordinate (B) at (2.2, 0);

  \fill[structurepurple] (N) circle (2.2pt);
  \node[above left=3pt, font=\small\bfseries, text=structurepurple] at (N) {$N$ (North Pole)};
  \fill[structurepurple] (A) circle (2.2pt) node[below left=2pt, font=\small, text=structurepurple] {$A$};
  \fill[structurepurple] (B) circle (2.2pt) node[above right=2pt, font=\small, text=structurepurple] {$B$};

  % Right angle indicators
  \draw[thin, gray!70] (0, -0.4) -- (0.25, -0.4) -- (0.25, -0.65);
  \draw[thin, gray!70] (1.95, 0) -- (1.95, 0.25) -- (2.2, 0.25);
  \draw[thin, gray!70] (0, 1.95) -- (0.25, 1.95) -- (0.25, 2.2);

  % Parallel transport vectors (trapred)
  % 1. Start vector at N pointing towards A (south)
  \draw[->, very thick, trapred] (N) -- ++(0, -1.0) node[left=2pt, font=\footnotesize] {$v_{\text{start}}$};

  % 2. Intermediate along NA
  \draw[->, thick, trapred] (0, 0.8) -- ++(0, -0.7);

  % 3. Vector at A pointing south
  \draw[->, very thick, trapred] (A) -- ++(0, -0.9) node[below=2pt, font=\footnotesize] {$v_A$};

  % 4. Intermediate along AB (perpendicular to equator, pointing south)
  \draw[->, thick, trapred] (1.2, -0.54) -- ++(0.3, -0.7);

  % 5. Vector at B pointing south (along negative tangent of meridian)
  \draw[->, very thick, trapred] (B) -- ++(0, -0.9) node[right=3pt, font=\footnotesize] {$v_B$};

  % 6. Intermediate along BN (pointing backwards along meridian)
  \draw[->, thick, trapred] (1.56, 1.56) -- ++(0.6, -0.6);

  % 7. Final vector at N pointing towards B (east, rotated 90 degrees!)
  \draw[->, very thick, trapred] (N) -- ++(1.0, 0) node[right=2pt, font=\footnotesize\bfseries] {$v_{\text{end}}$};

  % Holonomy rotation arc at N
  \draw[->, thick, structurepurple] (0, 1.5) arc[start angle=-90, end angle=0, radius=0.7] 
    node[midway, below right=1pt, font=\scriptsize] {$\Delta \theta = 90^\circ$};

  % Triangle label in interior
  \node[font=\footnotesize\bfseries, text=quotientorange!80!black] at (0.8, 0.7) {Area $= (\pi/2) R^2$};
\end{tikzpicture}
\caption{Parallel transport around a geodesic triangle on the sphere $S^2$. Carrying a vector along the path $N \to A \to B \to N$ without any in-surface twisting returns the vector rotated by $\Delta \theta = 90^\circ = \pi/2$ radians. This net rotation (holonomy) exactly matches the triangle's angle excess: $(90^\circ + 90^\circ + 90^\circ) - 180^\circ = 90^\circ$, which equals the total Gaussian curvature enclosed by the loop ($\iint_T K \, dA$).}\label{23:fig-parallel-transport-sphere}
\end{figure}

Consider the spherical triangle shown in Figure~\ref{23:fig-parallel-transport-sphere}:
\begin{enumerate}
\item Start at the North Pole $N$ with an arrow $v_{\text{start}}$ pointing South along the prime meridian toward $A$.
\item Walk down to the equator at $A$. Because you walk straight along the meridian, the arrow stays pointing South: $v_A$ points along the meridian past the equator.
\item Turn left $90^\circ$ and walk East along the equator to $B$. The arrow is perpendicular to your direction of travel. Parallel transport keeps it perpendicular to the equator at every step, so at $B$ it still points directly South ($v_B$).
\item Turn left $90^\circ$ and walk North along the meridian back to $N$. Your direction of motion is North, so the arrow pointing South points directly backward along your path.
\item When you arrive back at the North Pole $N$, the arrow $v_{\text{end}}$ points along the meridian $NB$---pointing East, rotated by exactly $90^\circ$ ($\pi/2$ radians) from your starting vector $v_{\text{start}}$!
\end{enumerate}

Even though the vector was parallel-transported without in-surface rotation along each segment, it returns rotated. This net angle of rotation around a closed loop is called \textbf{holonomy}.

Holonomy directly measures curvature. Gauss and Pierre Ossian Bonnet showed that the angle excess of any geodesic triangle $T$ with interior angles $\alpha_1, \alpha_2, \alpha_3$ equals the integral of the Gaussian curvature $K$ over the triangle's interior:
\[
\sum_{i=1}^3 \alpha_i - \pi = \iint_T K \, dA.
\]
For our triangle on a sphere of radius $R$:
\begin{itemize}
\item All three angles are right angles: $\alpha_1 = \alpha_2 = \alpha_3 = \pi/2$. The angle excess is $\frac{3\pi}{2} - \pi = \frac{\pi}{2}$.
\item The triangle covers one-eighth of the sphere, so its area is $\frac{1}{8}(4\pi R^2) = \frac{\pi}{2} R^2$.
\item The Gaussian curvature of a sphere of radius $R$ is $K = 1/R^2$.
\end{itemize}
Integrating the curvature yields:
\[
\iint_T K \, dA = \frac{1}{R^2} \left(\frac{\pi}{2} R^2\right) = \frac{\pi}{2}.
\]
The angle excess, the holonomy rotation, and the integrated curvature coincide. Figure~\ref{23:fig-three-geometries} compares how Gaussian curvature shapes geodesic triangles and the behavior of geodesics in positive, zero, and negative curvature.

\begin{figure}[htbp]
\centering
\begin{tikzpicture}[scale=0.95, >=Latex]
  % Column 1: Positive Curvature (x center = -4.3)
  \node[font=\small\bfseries, text=sourceblue, align=center] at (-4.3, 2.65) {Positive Curvature\\($K > 0$)};
  \node[font=\scriptsize\itshape, text=gray!80!black] at (-4.3, 2.05) {Spherical};

  % Spherical surface patch
  \draw[thick, sourceblue, fill=sourceblue!10] 
    (-5.5, 0.4) arc[start angle=180, end angle=0, x radius=1.2, y radius=1.2] -- cycle;
  \draw[dashed, sourceblue!50] 
    (-5.5, 0.4) arc[start angle=180, end angle=0, x radius=1.2, y radius=0.32];
  \draw[thick, sourceblue] 
    (-5.5, 0.4) arc[start angle=180, end angle=360, x radius=1.2, y radius=0.32];

  % Spherical triangle with bulging edges
  \fill[quotientorange!30] 
    (-4.3, 1.45) to[out=-105, in=55] (-5.0, 0.5) 
    to[out=10, in=170] (-3.6, 0.5) 
    to[out=125, in=-75] cycle;
  \draw[very thick, quotientorange!90!black] 
    (-4.3, 1.45) to[out=-105, in=55] (-5.0, 0.5) 
    to[out=10, in=170] (-3.6, 0.5) 
    to[out=125, in=-75] cycle;
  \node[font=\tiny\bfseries, text=quotientorange!90!black] at (-4.3, 0.8) {$\sum \alpha_i > \pi$};

  % Info card 1
  \node[draw=sourceblue!50, fill=sourceblue!5, rounded corners=3pt, inner sep=4pt, font=\scriptsize, align=center] at (-4.3, -0.85) {
    \textbf{Angle sum:} $\sum \alpha_i > 180^\circ$\\[2pt]
    \textbf{Geodesics:} Converge\\[2pt]
    \textbf{Archetype:} Sphere $S^2$
  };

  % Column 2: Zero Curvature (x center = 0.0)
  \node[font=\small\bfseries, text=targetgreen, align=center] at (0.0, 2.65) {Zero Curvature\\($K = 0$)};
  \node[font=\scriptsize\itshape, text=gray!80!black] at (0.0, 2.05) {Flat / Euclidean};

  % Flat plane patch
  \draw[thick, targetgreen, fill=targetgreen!10] 
    (-1.4, 0.2) -- (0.8, 0.2) -- (1.4, 1.6) -- (-0.8, 1.6) -- cycle;

  % Euclidean flat triangle
  \fill[quotientorange!30] 
    (0.0, 1.5) -- (-0.9, 0.35) -- (0.9, 0.35) -- cycle;
  \draw[very thick, quotientorange!90!black] 
    (0.0, 1.5) -- (-0.9, 0.35) -- (0.9, 0.35) -- cycle;
  \node[font=\tiny\bfseries, text=quotientorange!90!black] at (0.0, 0.65) {$\sum \alpha_i = \pi$};

  % Info card 2
  \node[draw=targetgreen!50, fill=targetgreen!5, rounded corners=3pt, inner sep=4pt, font=\scriptsize, align=center] at (0.0, -0.85) {
    \textbf{Angle sum:} $\sum \alpha_i = 180^\circ$\\[2pt]
    \textbf{Geodesics:} Parallel\\[2pt]
    \textbf{Archetype:} Flat Plane $\R^2$
  };

  % Column 3: Negative Curvature (x center = 4.3)
  \node[font=\small\bfseries, text=trapred, align=center] at (4.3, 2.65) {Negative Curvature\\($K < 0$)};
  \node[font=\scriptsize\itshape, text=gray!80!black] at (4.3, 2.05) {Hyperbolic / Saddle};

  % Saddle surface patch
  \draw[thick, trapred, fill=trapred!10] 
    (3.0, 1.5) to[out=-20, in=200] (5.6, 1.5) 
    to[out=-80, in=60] (5.4, 0.2) 
    to[out=170, in=10] (3.2, 0.2) 
    to[out=120, in=-100] cycle;
  \draw[thin, dashed, trapred!60] 
    (3.0, 1.5) to[out=-40, in=160] (4.3, 0.6) to[out=-20, in=180] (5.4, 0.2);

  % Hyperbolic triangle with pinched edges
  \fill[quotientorange!30] 
    (4.3, 1.35) to[out=-80, in=30] (3.5, 0.35) 
    to[out=10, in=170] (5.1, 0.35) 
    to[out=150, in=-100] cycle;
  \draw[very thick, quotientorange!90!black] 
    (4.3, 1.35) to[out=-80, in=30] (3.5, 0.35) 
    to[out=10, in=170] (5.1, 0.35) 
    to[out=150, in=-100] cycle;
  \node[font=\tiny\bfseries, text=quotientorange!90!black] at (4.3, 0.7) {$\sum \alpha_i < \pi$};

  % Info card 3
  \node[draw=trapred!50, fill=trapred!5, rounded corners=3pt, inner sep=4pt, font=\scriptsize, align=center] at (4.3, -0.85) {
    \textbf{Angle sum:} $\sum \alpha_i < 180^\circ$\\[2pt]
    \textbf{Geodesics:} Diverge\\[2pt]
    \textbf{Archetype:} Hyperbolic $\mathbb{H}^2$
  };
\end{tikzpicture}
\caption{The geometry of surfaces with positive, zero, and negative Gaussian curvature $K$. In positive curvature (left, sphere), geodesics converge and triangle angles sum to more than $180^\circ$, with the excess measuring enclosed curvature. In zero curvature (center, plane), Euclid's parallel postulate holds and triangle angles sum to exactly $180^\circ$. In negative curvature (right, saddle), geodesics diverge and triangle angles sum to strictly less than $180^\circ$, with the deficit proportional to area.}\label{23:fig-three-geometries}
\end{figure}

\begin{intuitionbox}
\textbf{Pause and think.} Suppose you walk around a small circle of radius $r$ on a flat parking lot, holding a compass that points steadily North. When you return to your starting point, does the compass needle point North?

Yes. On flat Euclidean ground, parallel-transporting a direction vector around any closed loop returns it with zero rotation. But on a curved sphere of radius $R$, the compass needle precesses: the rotation angle is proportional to the enclosed area divided by $R^2$. As the circle shrinks toward zero radius ($r \to 0$), the enclosed area vanishes and the rotation angle approaches zero, recovering the local flatness of the tangent plane.
\end{intuitionbox}

\begin{historybox}
Before Gauss, curvature was generally treated extrinsically: a curve or surface was thought to bend only because it was immersed in an ambient higher-dimensional space, like a circle in the plane or a twisted ribbon in three-dimensional space. Gauss's 1827 \emph{Theorema Egregium} proved that curvature is intrinsic, computable entirely from measurements along the surface itself without reference to an external container.

On June 10, 1854, Gauss sat in the audience at the University of G\"ottingen to evaluate the Habilitation lecture of his former student, Bernhard Riemann. Riemann's lecture, titled \emph{\"Uber die Hypothesen, welche der Geometrie zu Grunde liegen} (``On the Hypotheses Which Lie at the Bases of Geometry''), marked a turning point in geometry.

At the time, the dominant philosophy of Immanuel Kant maintained that Euclidean 3-dimensional space was an immutable, \emph{a priori} necessary framework of human cognition---a truth known prior to all experience. Riemann challenged this view. He showed that an $n$-dimensional space (a \emph{manifold}) requires no external container or higher dimensions. It is defined intrinsically, endowed with a metric tensor measuring infinitesimal distances between nearby points. In higher dimensions, curvature is described by the Riemann curvature tensor ($R^i_{jk\ell}$). Riemann argued that whether physical space is flat or curved, three-dimensional or higher-dimensional, is an empirical question to be settled by physical experiment.

Six decades later, Albert Einstein---assisted by Marcel Grossmann and Hermann Minkowski---found that Riemann's geometry was the mathematical language of general relativity. Spacetime is a four-dimensional pseudo-Riemannian manifold (where time enters the metric with an opposite sign relative to space). Matter and energy dictate the curvature of spacetime via Einstein's field equations, and that curvature dictates how matter moves: gravitation is described as the intrinsic curvature of spacetime itself.

The geometry initiated by Gauss and Riemann remains an active area of research. In 1969, Roger Penrose formulated the \textbf{cosmic censorship conjecture}: are physical gravitational singularities (points where curvature invariants diverge) always hidden behind black hole event horizons, or can ``naked singularities'' be exposed to the outside universe? In topology, while Grigori Perelman proved the 3-dimensional Poincar\'e conjecture in 2003 using Richard Hamilton's Ricci curvature flow, the \textbf{smooth 4-dimensional Poincar\'e conjecture} remains open: can a four-dimensional manifold be homeomorphic to a 4-sphere while carrying an incompatible smooth differential structure? Finally, uniting the smooth geometric spacetime of general relativity with quantum mechanics remains an unsolved challenge in theoretical physics.
\end{historybox}

\subsection{Differential geometry---and singular spaces}

Differential geometry lets us write calculus and physics in a form that does not depend on an observer's choice of coordinates. Differential forms and tensor fields express physical laws invariantly on curved spaces, and multivariable integral theorems---such as Stokes' theorem and the divergence theorem---carry over from $\R^n$ to smooth manifolds, where they govern fluid flow, electromagnetism, and gravity.

That machinery, however, depends on smoothness everywhere. Every small patch must look smoothly like flat Euclidean space, with differentiable transition maps on overlaps. As soon as a shape develops a sharp pinch or crosses itself---like the nodal curve $y^2 = x^3 + x^2$ at the origin---there is no single flat tangent space to linearize around, and smooth manifold calculus breaks down at the singular point. Nor do derivatives directly track discrete algebraic features such as repeated intersection multiplicities.

To study geometric sets defined by polynomial equations that may include self-intersections, cusps, and singularities, Chapter~24 introduces \textbf{Algebraic Geometry Preview: Equations as Shapes}, replacing calculus with polynomial rings, vanishing loci, and ideals. Later, Chapter~38 develops this perspective in \textbf{Second Pass: Geometry, Coordinates, and What Is Real}.

\subsection{Exercises}

\begin{enumerate}
\item \textbf{Geodesics on a cylinder.} (Mental check)
An ant crawls on the surface of an infinite cylinder $x^2 + y^2 = 1$ in $\R^3$.
\begin{enumerate}
\item If you cut the cylinder along the vertical line $x = 1, y = 0$ and unroll it flat onto a desk, straight lines on the flat paper correspond to which curves on the cylinder?
\item An ant wants to walk from $(1, 0, 0)$ to $(1, 0, 4\pi)$ along a geodesic path. How many complete loops around the cylinder might such a geodesic make, and how does its total length depend on the number of turns $k \in \Z$?
\item Explain why the Gaussian curvature of the cylinder is $K = 0$ everywhere, even though the cylinder appears bent to an outside observer in $\R^3$.
\end{enumerate}

\item \textbf{Spherical excess and area.} (Routine practice)
On the unit sphere $S^2$ (radius $R = 1$, Gaussian curvature $K = 1$), a surveyor measures a geodesic triangle whose interior angles are $\alpha = 75^\circ$, $\beta = 85^\circ$, and $\gamma = 80^\circ$.
\begin{enumerate}
\item Convert these angles to radians and compute the spherical excess $E = \alpha + \beta + \gamma - \pi$.
\item Using the Gauss-Bonnet formula $E = \iint_T K \, dA$, find the exact surface area of the triangle.
\item What fraction of the sphere's total surface area ($4\pi$) does this triangle cover?
\end{enumerate}

\item \textbf{The circle atlas.} (Routine practice)
Consider the unit circle $S^1 = \{(x, y) \in \R^2 : x^2 + y^2 = 1\}$. Let $N = (0, 1)$ be the North pole and $S = (0, -1)$ be the South pole. Stereographic projection from $N$ onto the $x$-axis gives the chart map $\phi_1: S^1 \setminus \{N\} \to \R$ with $\phi_1(x, y) = u = \frac{x}{1 - y}$. Stereographic projection from $S$ gives $\phi_2: S^1 \setminus \{S\} \to \R$ with $\phi_2(x, y) = v = \frac{x}{1 + y}$.
\begin{enumerate}
\item Show that on the overlap $(S^1 \setminus \{N\}) \cap (S^1 \setminus \{S\})$, the coordinates satisfy $u \cdot v = 1$, so the transition map is $v = \tau_{12}(u) = 1/u$.
\item Show that the transition map is infinitely differentiable ($C^\infty$) on its overlap domain $\R \setminus \{0\}$.
\item Conclude that this two-chart atlas gives $S^1$ the structure of a 1-dimensional smooth manifold.
\end{enumerate}

\item \textbf{The metric tensor and lengths on the sphere.} (Computation)
In spherical coordinates $(\theta, \phi)$ with colatitude $\theta \in (0, \pi)$ and longitude $\phi \in [0, 2\pi)$, the line element on a sphere of radius $R$ is:
\[
ds^2 = R^2 \, d\theta^2 + R^2 \sin^2\theta \, d\phi^2.
\]
\begin{enumerate}
\item Write the matrix of metric coefficients $[g_{ij}]$ in the $(\theta, \phi)$ coordinate basis.
\item Consider a path along a line of constant latitude $\theta = \theta_0$: $\gamma(t) = (\theta_0, t)$ for $t \in [0, 2\pi]$. Compute the velocity vector $\gamma'(t)$ and the speed $\sqrt{g(\gamma'(t), \gamma'(t))}$.
\item Integrate the speed from $t = 0$ to $2\pi$ to find the circumference of the latitude circle as a function of $\theta_0$. For which value of $\theta_0$ is the circumference maximized, and what is its value?
\end{enumerate}

\item \textbf{Parallel transport on a cone.} (Deeper challenge)
Take a flat sheet of paper, cut out a circular sector of angle $\Delta\phi \in (0, 2\pi)$, and tape the cut edges together without stretching to form a cone with apex angle $\alpha$.
\begin{enumerate}
\item What is the Gaussian curvature $K$ at any point on the cone away from the sharp apex? Explain your reasoning using Gauss's \emph{Theorema Egregium}.
\item A surveyor carries a tangent vector along a simple closed loop that circles once around the apex of the cone, parallel-transporting it along the path. By what angle has the vector rotated when it returns to its starting point? (Hint: unroll the cone back onto the flat paper).
\item How does this result show that all of the cone's curvature is concentrated at the single singular point at the apex?
\end{enumerate}
\end{enumerate}

\newpage
